\section*{Chapitre \ref{ch:g4:what-a-fire-needs} --- La combustion~: ce dont un feu a besoin}

\begin{solution}{exo:g4:what-a-fire-needs:1}
Le \omterm{def:g4:what-a-fire-needs:fuel}{combustible}, l'oxygène (de l'\omterm{def:g4:air-a-mixture-of-gases:air}{air}) et la chaleur (pour démarrer le
feu).
\end{solution}

\begin{solution}{exo:g4:what-a-fire-needs:2}
Le bois, la cire de bougie, l'essence et le papier sont des
\omterm{def:g4:what-a-fire-needs:fuel}{combustibles}. Le sable, le verre et l'eau ne brûlent pas.
\end{solution}

\begin{solution}{exo:g4:what-a-fire-needs:3}
L'oxygène~: la couverture empêche l'\omterm{def:g4:air-a-mixture-of-gases:air}{air} d'atteindre le feu.
\end{solution}

\begin{solution}{exo:g4:what-a-fire-needs:4}
La chaleur~: l'eau refroidit tant le bois en feu qu'il ne peut plus
\omterm{def:g4:what-a-fire-needs:burning}{brûler}.
\end{solution}

\begin{solution}{exo:g4:what-a-fire-needs:5}
Trois parmi~: la cendre, la fumée, le dioxyde de carbone, la vapeur
d'eau.
\end{solution}

\begin{solution}{exo:g4:what-a-fire-needs:6}
Souffler apporte plus d'\omterm{def:g4:air-a-mixture-of-gases:air}{air}, donc plus d'oxygène, au \omterm{def:g4:what-a-fire-needs:fuel}{combustible} chaud,
et il se remet à \omterm{def:g4:what-a-fire-needs:burning}{brûler}.
\end{solution}

\begin{solution}{exo:g4:what-a-fire-needs:7}
De l'eau~: la cire qui \omterm{def:g4:what-a-fire-needs:burning}{brûle} fabrique de la vapeur d'eau, qui se
transforme en gouttelettes sur le verre froid.
\end{solution}

\begin{solution}{exo:g4:what-a-fire-needs:8}
Sortir, en restant baissé sous la fumée et en fermant les portes
derrière soi, puis appeler les pompiers une fois dehors.
\end{solution}

\begin{solution}{exo:g4:what-a-fire-needs:9}
Le côté \omterm{def:g4:what-a-fire-needs:fuel}{combustible}~: une fois le gaz fermé, le feu n'a plus rien à
\omterm{def:g4:what-a-fire-needs:burning}{brûler}.
\end{solution}

\begin{solution}{exo:g4:what-a-fire-needs:10}
Le \omterm{def:g4:what-a-fire-needs:fuel}{combustible}~: comme il n'y a plus d'arbres dans la bande, le feu
manque de \omterm{def:g4:what-a-fire-needs:fuel}{combustible} quand il l'atteint.
\end{solution}

\begin{solution}{exo:g4:what-a-fire-needs:11}
Non. La plus grande partie du bois, avec l'oxygène de l'\omterm{def:g4:air-a-mixture-of-gases:air}{air}, est
devenue de nouvelles substances qui sont des gaz --- du dioxyde de
carbone, de la vapeur d'eau et de la fumée ---, et elles sont parties
dans l'\omterm{def:g4:air-a-mixture-of-gases:air}{air}. Seule la cendre reste.
\end{solution}
